% © 2025 Ing. David Jaroš — CC BY-NC-ND 4.0
%
% This work is licensed under a Creative Commons Attribution-NonCommercial-NoDerivatives
% 4.0 International License (CC BY-NC-ND 4.0).
%
% B Coefficient from Loop Expansion in UBT — Candidate 2: Non-Commutative Correction
% Track: CORE — α derivation
% Date: 2026-03-04

\ifdefined\INCLUDEMODE
  % Being included in another document — skip preamble
\else
  \documentclass[11pt,a4paper]{article}
  \usepackage{amsmath,amssymb,amsthm}
  \usepackage{hyperref}

  \newtheorem{theorem}{Theorem}
  \newtheorem{proposition}{Proposition}
  \newtheorem{lemma}{Lemma}
  \theoremstyle{remark}
  \newtheorem{remark}{Remark}

  \title{Candidate 2: Non-Commutative Correction from
    \texorpdfstring{$[D_\mu, D_\nu] \neq 0$}{[D,D] != 0}\\
    \large B Coefficient from Loop Expansion — Biquaternionic Curvature Analysis}
  \author{Ing.\ David Jaroš}
  \date{2026-03-04}

  \begin{document}
  \maketitle
\fi

\section{Candidate 2: Non-Commutative Correction from
  \texorpdfstring{$[D_\mu,D_\nu]\neq 0$}{[D,D] != 0}}
\label{sec:candidate2_noncommutative}

\subsection{Physical Motivation}

In standard QED on flat Minkowski space, the covariant derivatives satisfy:
\begin{equation}
  [D_\mu, D_\nu] = i e F_{\mu\nu},
  \label{eq:QED_commutator}
\end{equation}
where $F_{\mu\nu} = \partial_\mu A_\nu - \partial_\nu A_\mu$ is the electromagnetic field
strength. This is the \emph{only} contribution: the base space is flat and commutative.

In UBT on $\mathbb{C} \otimes \mathbb{H}$, the biquaternionic structure introduces an additional
curvature contribution:
\begin{equation}
  [D_\mu, D_\nu] = i e F_{\mu\nu} + \Omega_{\mu\nu},
  \label{eq:UBT_commutator}
\end{equation}
where $\Omega_{\mu\nu}$ encodes the biquaternionic curvature absent in QED. In the QED limit
($\psi \to 0$, imaginary sector collapses), $\Omega_{\mu\nu} \to 0$ and standard QED is recovered.

\subsection{Biquaternionic Covariant Derivative in UBT}

The UBT covariant derivative on $\mathbb{C} \otimes \mathbb{H}$ is:
\begin{equation}
  D_\mu = \partial_\mu + i e A_\mu + \Gamma_\mu^{\mathbb{H}},
  \label{eq:UBT_D}
\end{equation}
where $\Gamma_\mu^{\mathbb{H}}$ is the biquaternionic connection arising from the non-trivial
imaginary metric $h_{\mu\nu} = \mathrm{Im}[\mathcal{G}_{\mu\nu}]$.

For the two-mode winding vacuum ansatz
(canonical/geometry/biquaternionic\_vacuum\_solutions.tex \S1.3):
\begin{equation}
  \Theta(\psi) = \Theta_0\,e^{i\psi/R_\psi} + \Theta_1\,e^{2i\psi/R_\psi},
  \label{eq:two_mode_ansatz}
\end{equation}
the imaginary metric component is (see canonical/geometry/biquaternionic\_vacuum\_solutions.tex \S1.3):
\begin{equation}
  h_{\psi\psi}(\psi) = \frac{2}{R_\psi^2}\,\sin\!\Bigl(\frac{\psi}{R_\psi}\Bigr)\cdot
    \mathrm{Im}\bigl[\mathrm{Sc}(\Theta_0 \Theta_1^\dagger)\bigr].
  \label{eq:h_psi_psi}
\end{equation}
For the canonical example ($\mathrm{Im}[\mathrm{Sc}(\Theta_0\Theta_1^\dagger)] = 1$,
$R_\psi = 1$): $h_{\psi\psi}(\psi) = 2\sin\psi$.

\subsection{Computation of \texorpdfstring{$\Omega_{\mu\nu}$}{Omega\_munu}}

The biquaternionic curvature two-form $\Omega_{\mu\nu}$ arises from the non-commutativity of
the internal $\mathbb{H}$ algebra. For the gauge-singlet sector (electromagnetic $\mathrm{U}(1)$),
the relevant contribution is the Riemann curvature from the imaginary metric $h_{\mu\nu}$:
\begin{equation}
  \Omega_{\mu\nu} = R^\psi{}_{\mu\nu\psi}(h)\,d\psi,
  \label{eq:Omega_def}
\end{equation}
where $R^\psi{}_{\mu\nu\psi}(h)$ is the Riemann tensor computed from $h_{\mu\nu}$ as a
perturbation of flat Minkowski space.

\textbf{Step 1: Compute $\Omega_{\mu\nu}$ in the UBT vacuum.}

For the imaginary metric component~\eqref{eq:h_psi_psi}, the only non-trivial component in
the $(\mu,\nu) = (\psi,\psi)$ sector is:
\begin{equation}
  \partial_\psi^2 h_{\psi\psi} = -\frac{2}{R_\psi^4}\,\sin\!\Bigl(\frac{\psi}{R_\psi}\Bigr)
    \cdot \mathrm{Im}\bigl[\mathrm{Sc}(\Theta_0\Theta_1^\dagger)\bigr].
  \label{eq:d2h}
\end{equation}

The effective curvature in the 4D sector comes from integrating over $\psi$ (Kaluza--Klein
dimensional reduction):
\begin{equation}
  \bar{\Omega}_{\mu\nu}^{\mathrm{eff}} = \frac{1}{2\pi R_\psi}\int_0^{2\pi R_\psi}
    \Omega_{\mu\nu}(\psi)\,d\psi.
  \label{eq:Omega_effective}
\end{equation}
For the sinusoidal $h_{\psi\psi}$, the average $\langle \Omega_{\mu\nu} \rangle = 0$
(zero mode vanishes by symmetry: $\int_0^{2\pi} \sin\psi\,d\psi = 0$).

This means: \emph{the zero-mode contribution to $\Omega_{\mu\nu}^{\mathrm{eff}}$ vanishes.}

\textbf{Step 2: UV-divergent coefficient $\delta B_{\mathrm{NC}}$.}

The non-commutative correction to the vacuum polarisation has the form:
\begin{equation}
  \delta \Pi_{\mathrm{NC}}(k^2) = C_{\mathrm{NC}} \cdot
    \frac{\mathrm{Tr}\bigl([\Omega_{\mu\nu}]^2\bigr)}{k^2},
  \label{eq:delta_Pi_NC}
\end{equation}
where the trace runs over biquaternionic indices and the loop momentum.

The UV-divergent piece of $\delta\Pi_{\mathrm{NC}}$ determines $\delta B_{\mathrm{NC}}$:
\begin{equation}
  \delta B_{\mathrm{NC}} = \lim_{\epsilon\to 0} \Bigl[
    \frac{\delta\Pi_{\mathrm{NC}}(k^2)}{k^2}\Bigr]_{\mathrm{UV}}.
  \label{eq:delta_B_NC}
\end{equation}

\textbf{Step 3: Evaluate $\mathrm{Tr}(\Omega_{\mu\nu}^2)$.}

In the $\psi$-direction sector, the relevant trace is:
\begin{equation}
  \mathrm{Tr}\bigl(\Omega_{\psi\psi}^2\bigr) = \int_0^{2\pi R_\psi}
    [h_{\psi\psi}(\psi)]^2\,\frac{d\psi}{2\pi R_\psi}
  = \bigl[\mathrm{Im}(\mathrm{Sc}(\Theta_0\Theta_1^\dagger))\bigr]^2 \cdot
    \frac{4}{R_\psi^4} \cdot \frac{1}{2},
  \label{eq:Tr_Omega_sq}
\end{equation}
since $\frac{1}{2\pi}\int_0^{2\pi}\sin^2\psi\,d\psi = 1/2$.

For the canonical vacuum ($\mathrm{Im}[\mathrm{Sc}(\Theta_0\Theta_1^\dagger)] = 1$, $R_\psi = 1$):
\begin{equation}
  \mathrm{Tr}\bigl(\Omega_{\psi\psi}^2\bigr) = 2.
  \label{eq:Tr_Omega_sq_canonical}
\end{equation}

The UV cutoff in UBT is set geometrically by $\Lambda = 1/R_\psi$. The contribution to the
vacuum polarisation from $\Omega$ at scale $\Lambda$ is:
\begin{equation}
  \delta \Pi_{\mathrm{NC}} \sim C_{\mathrm{NC}} \cdot
    \frac{\mathrm{Tr}(\Omega_{\psi\psi}^2)}{(1/R_\psi)^2}
  = C_{\mathrm{NC}} \cdot 2 \cdot R_\psi^2 = C_{\mathrm{NC}} \cdot 2
  \quad (R_\psi = 1),
  \label{eq:delta_Pi_NC_cutoff}
\end{equation}
giving
\begin{equation}
  \delta B_{\mathrm{NC}} = C_{\mathrm{NC}} \cdot 2 \cdot B_0.
  \label{eq:delta_B_NC_result}
\end{equation}

\subsection{Checking Against the Required Value}

For $B_0 + \delta B_{\mathrm{NC}} = 46.3$:
\begin{equation}
  B_0(1 + 2 C_{\mathrm{NC}}) = 46.3
  \implies 2 C_{\mathrm{NC}} = \frac{46.3}{B_0} - 1 = 1.844 - 1 = 0.844
  \implies C_{\mathrm{NC}} = 0.422.
  \label{eq:C_NC_required}
\end{equation}
A coefficient $C_{\mathrm{NC}} \approx 0.42$ is geometrically natural if it arises from a
loop integral over the biquaternionic algebra. For comparison, the standard QED two-loop
coefficient in the Schwinger-Dyson expansion is $O(1)$.

\subsection{Explicit Derivation Requirement}

The coefficient $C_{\mathrm{NC}}$ must be computed from the one-loop integral with the
non-commutative propagator. The relevant Feynman diagram has the structure:
\begin{equation}
  C_{\mathrm{NC}} = \frac{1}{(4\pi)^2}\int_0^1 dx\,2x(1-x)\,
    \mathrm{Tr}_{\mathbb{H}}\bigl[\gamma^\mu\Gamma_\mu^{\mathbb{H}}\bigr]
    + \text{(biquaternionic index contractions)},
  \label{eq:C_NC_integral}
\end{equation}
where $\mathrm{Tr}_{\mathbb{H}}$ denotes the trace over quaternionic indices.

Computing $\mathrm{Tr}_{\mathbb{H}}[\gamma^\mu \Gamma_\mu^{\mathbb{H}}]$ requires:
\begin{enumerate}
  \item Writing $\Gamma_\mu^{\mathbb{H}}$ explicitly from the two-mode vacuum
    connection formula (canonical/geometry, \S2).
  \item Evaluating the trace using the biquaternionic inner product.
  \item Extracting the UV-divergent piece via dimensional regularisation.
\end{enumerate}
This computation is not yet completed in the present document.

\subsection{QED Limit Verification}

In the QED limit ($\psi \to 0$):
\begin{itemize}
  \item $\Theta_0, \Theta_1 \to$ real-valued quaternions.
  \item $\mathrm{Im}[\mathrm{Sc}(\Theta_0 \Theta_1^\dagger)] \to 0$.
  \item $h_{\psi\psi} \to 0$.
  \item $\Omega_{\mu\nu} \to 0$.
  \item $\delta B_{\mathrm{NC}} \to 0$.
\end{itemize}
Therefore:
\begin{equation}
  B\bigl(N_{\mathrm{eff}} = 1, \delta_{\mathrm{NC}} = 0\bigr)
  = B_0\bigl(N_{\mathrm{eff}} = 1\bigr) = \frac{2\pi}{3}.
  \label{eq:qed_limit_C2}
\end{equation}
The QED anchor is satisfied. \checkmark

\subsection{Status Assessment}

\begin{proposition}[NC correction: status]
\label{prop:NC_status}
The non-commutative correction $\delta B_{\mathrm{NC}}$ from $[D_\mu, D_\nu] \neq 0$ satisfies:
\begin{enumerate}
  \item The zero-mode contribution vanishes ($\langle \Omega_{\mu\nu} \rangle = 0$).
  \item The quadratic contribution $\mathrm{Tr}(\Omega^2) = 2$ (canonical vacuum, $R_\psi = 1$).
  \item The required $C_{\mathrm{NC}} = 0.422$ is geometrically plausible but \emph{not yet derived}
    from the biquaternionic loop integral.
  \item The formula $\delta B_{\mathrm{NC}} = C_{\mathrm{NC}} \cdot \mathrm{Tr}(\Omega^2) \cdot B_0$
    vanishes in the QED limit.
\end{enumerate}
\end{proposition}

\textbf{Conclusion.} Candidate 2 is \textbf{not ruled out}. The non-commutative correction
could in principle close the gap if $C_{\mathrm{NC}} \approx 0.422$. However, this requires an
explicit computation of the biquaternionic loop integral \eqref{eq:C_NC_integral} with
connection $\Gamma_\mu^{\mathbb{H}}$ from the two-mode vacuum. This computation is deferred
to a dedicated calculation.

\begin{center}
\textbf{Candidate 2 (NC correction): STATUS — REQUIRES EXPLICIT LOOP INTEGRAL}

$C_{\mathrm{NC}} = 0.422$ required. Geometrically plausible. Not yet derived.
\end{center}

\ifdefined\INCLUDEMODE\else
\end{document}
\fi
